\chapter{Futures Market Making}\label{ch:hf:futures-market-making}

In a short-rate futures contract the queue at the best price can hold tens of thousands of lots, and fills are shared pro rata. A market maker who
wants a hundred lots shows more, because its share of each fill is its size over everyone's, and every other maker reasons the same way. In this
chapter's game, ten makers who each want 100 lots show 230; fifteen show 620; twenty show as much as the exchange allows. Each is worse off than if
all had shown what they wanted, and the depth at the best price reaches 245 times the average order, as Field and Large found in real pro-rata
markets.

\section{Pro-rata books and size}

Book 1, chapter 19 sets out the allocation rules: time priority (first in, first filled), pro rata (each resting order gets its share of the
incoming size), and the hybrids that give a top order, the first to improve the price, or a lead market maker a slice first. Under time priority a
market maker's decision is where to stand in the queue (chapter 5); under pro rata it is how much to show.

\begin{dated}{2026-09}{How one exchange allocates}\label{dat:hf:futures-market-making:cme}
As a market-data vendor's guide to CME's matching algorithms described them in August 2025, pro-rata allocation matches an incoming aggressor ``based
on each resting order's pro-rated percentage''; allocations are rounded down to the nearest integer, including zero, and excess lots are allocated
first in, first out. In the algorithms that include it, the first order that improves the market (the top order) is matched first, followed by any
lead market maker allocation. The guide listed pro rata for some calendar-spread books and a configurable algorithm for outrights such as the
two-year Treasury note future.
\end{dated}

The allocation of an aggressor of 250 lots against four resting orders of 100 (the top order), 400, 300 and 200 shows the difference: time priority
gives the first two 100 and 150; pure pro rata gives 25, 100, 75 and 50; top order then pro rata gives 100, 67, 50 and 33 (\texttt{firm.match}).

\begin{definition}[Over-quoting]\label{def:hf:futures-market-making:over}
\emph{Over-quoting}\index{over-quoting} is showing more size at a price than one wishes to trade there, in a book that allocates fills in
proportion to size, in order to receive a larger share of each fill; it is profitable for each maker alone and self-defeating for all together.
\end{definition}

The game (\cref{lst:hf:futures-market-making:game}): $n$ makers rest at the best price; an aggressor of $X$ lots arrives, lognormal with a median of
200; maker $i$, showing $s_i$ of a total $S$, receives $s_i\min(X,S)/S$. It earns one tick a lot and pays $\gamma/2=0.005$ per squared lot of
inventory. Alone, or first in a time-priority queue, it would show 100 lots, the size whose marginal fill is worth its marginal risk. Against others,
showing more buys a larger share of every order; the symmetric equilibrium is found by best responses.

\begin{omfigure}
\begin{tikzpicture}
\begin{axis}[width=11cm,height=5.2cm,xlabel={\small makers at the best price},ylabel={\small size shown / size wanted},ymode=log,log basis y=10,
  xmin=0,xmax=21,ymin=0.8,ymax=130,ytick={1,2,5,10,20,50},yticklabels={1,2,5,10,20,50},tick label style={font=\footnotesize},grid=major,
  grid style={gray!25},table/col sep=comma]
\addplot[omThm,thick,mark=*,mark size=1.4pt] table[x=n,y=factor]{figdata/hft/15-futures-market-making/overquote.csv};
\draw[gray,dashed] (axis cs:0,50) -- (axis cs:21,50);
\node[font=\scriptsize,anchor=south] at (axis cs:10,50) {exchange's maximum order};
\end{axis}
\end{tikzpicture}

\omcaption{The \omterm{def:hf:futures-market-making:over}{over-quoting} factor, the equilibrium size each maker shows divided by the 100 lots it would show under time priority, against the
number of makers at the best price (log scale); aggressors lognormal with a median of 200 lots, one tick earned a lot, an \omterm{def:hf:inventory-models:penalty}{inventory penalty} $\gamma=0.01$
(0.005 per squared lot), sizes capped at 5,000. Data: \texttt{hf\_futures.game}.}\label{fig:hf:futures-market-making:over}
\end{omfigure}

The factor rises from 1.1 with two makers to 2.3 with ten and 6.2 with fifteen; with twenty, the spiral stops only at the cap
(\cref{fig:hf:futures-market-making:over}). Every maker's expected utility is lower than if all had shown 100: 19.5 against 22.3 with ten makers.
Field and Large found exactly the symptoms in four one-tick pro-rata futures markets: depth at the quotes exceeding the mean market order by two orders
of magnitude, and cancellation rates above 96\% of the lots offered. They modelled the cause as the strategic complementarity in the game: each
trader risks overtrading with an oversized order that it expects mostly to cancel.

The risk is the rare large order. An aggressor at the 99.9th percentile of the simulated sizes (7,907 lots) fills each of ten makers with 230 lots,
2.3 times what each wanted, all at once and all on the same side: the whole market becomes long together.

\section{Calendar spreads and implied prices}

A calendar spread (Book 1, chapter 19) buys one expiry and sells another. Exchanges list spread books alongside the outrights and link them through
implied prices: a spread's implied-in price comes from the two outright books (bid of the front minus offer of the back), an outright's implied-out
price from a spread book and the other outright. The exchange's matching engine can trade through these links, so an order in one book may fill
against orders in two others.

\section{Arbitrage between spread and outright books}

\begin{definition}[Implied-price arbitrage]\label{def:hf:futures-market-making:implied}
\emph{Implied-price arbitrage}\index{implied-price arbitrage} is trading a spread book against its two outright books when the spread's direct
price is outside the range implied by the outrights (or the reverse), before the exchange's implied matching or other traders close the gap.
\end{definition}

Where the exchange computes implied prices for every link, many such gaps are closed by the engine itself; where it does not (for some spreads,
across exchanges, or between a spread and a pack of several outrights), they are closed by traders. The gaps come from lags: the spread book's market
makers quote around the outrights' difference as they last saw it. In the chapter's strip of four quarterly futures, outright books one tick wide and
spread books one tick wide around a lagged difference, arbitrages appear when the lag spans enough market-data events for the curve's slope to move
two ticks: none with a lag of one event, 0.9 per thousand spread-events with five, 5.9 with ten and 50 with thirty, each worth about a tick
(\cref{fig:hf:futures-market-making:implied}).

\begin{omfigure}
\begin{tikzpicture}
\begin{axis}[width=11cm,height=4.8cm,xlabel={\small spread books' lag (market-data events)},ylabel={\small arbitrages per 1,000},xmin=0,xmax=31,
  ymin=0,ymax=55,tick label style={font=\footnotesize},grid=major,grid style={gray!25},table/col sep=comma]
\addplot[omThm,thick,mark=*,mark size=1.4pt] table[x=lag,y=per1000]{figdata/hft/15-futures-market-making/implied.csv};
\end{axis}
\end{tikzpicture}

\omcaption{\omterm{def:hf:futures-market-making:implied}{Implied-price arbitrages} between three calendar-spread books and four outright books, per thousand spread-book observations, against how
many market-data events the spread books' quotes lag the outrights; the curve's level and slope move 0.6 and 0.4 ticks an event; 5,000 events. Data:
\texttt{hf\_futures.implied}.}\label{fig:hf:futures-market-making:implied}
\end{omfigure}

\section{Rates futures strips}

A short-rate futures strip (One Quant Book 2, chapter 8) is a curve of quarterly contracts, each a forward rate for a three-month period, traded
as outrights, calendar spreads, butterflies, packs and bundles. A market maker in the strip quotes all of them from one model of the curve (a few
factors: level, slope, curvature, plus the meeting dates of the central bank) and hedges any fill in whichever instrument is cheapest. The strip's
books are pro rata or hybrid in several markets, so the \omterm{def:hf:futures-market-making:over}{over-quoting} of the previous sections applies to each; the model of the curve is what keeps
the maker's quotes in hundreds of books consistent with each other.

\section{Inter-commodity spreads}

Two related contracts (two points on a yield curve in different contracts, crude oil and its products, corn and ethanol) trade as an
inter-commodity spread when their prices share a factor. The exchange's margin credit for the spread (Book 1, chapter 20) lowers its cost of
carry. A market maker quotes the spread from its model of the relation and hedges each fill in the leg that is cheaper to trade; the risk is the
relation's own drift, which no hedge removes.

\section{Strategy files}

\begin{strategyfile}{Pro-rata outright quoting}\label{strat:hf:futures-market-making:prorata}
\sfield{Who pays you, and why} Hedgers and directional traders who cross a one-tick spread in size.
\sfield{Instruments and venues} Short-rate and bond futures outrights on pro-rata or hybrid books.
\sfield{Signal} The fair value inside the tick (chapter 2) and the probability of a large aggressor, from order flow and the other books.
\sfield{Sizing and execution} Show the equilibrium size, not the wanted size; cancel ahead of predicted large orders; seek top-order and lead-market-maker
priority where available.
\sfield{Costs} Exchange fees per lot, the inventory from large fills, message limits from cancelling most of what is shown.
\sfield{How it dies} The rare very large order: ten makers each filled 2.3 times their wanted size by one aggressor.
\sfield{Horizon, capacity, infrastructure} Seconds to minutes; capacity large; a curve model and fast cancellation.
\sfield{Backtest honestly} The allocation algorithm and its parameters by product; others' sizes (unobservable) modelled, not assumed away.
\sfield{Sources} Field and Large (2008); the dated box.
\end{strategyfile}

\begin{strategyfile}{Calendar-spread quoting with implied legs}\label{strat:hf:futures-market-making:spread}
\sfield{Who pays you, and why} Hedgers rolling positions from one expiry to the next, who trade the spread to avoid legging.
\sfield{Instruments and venues} The exchange's calendar-spread books and their outrights.
\sfield{Signal} The outrights' difference from the curve model, with the exchange's implied prices as the competing quote.
\sfield{Sizing and execution} Quote the spread inside the implied market; hedge a fill in the outrights when the model says the spread will not
mean-revert quickly.
\sfield{Costs} Fees on the spread and on any hedge; queue position behind implied orders.
\sfield{How it dies} Roll periods concentrate flow and competitors; the implied engine undercuts a slow quote.
\sfield{Horizon, capacity, infrastructure} Seconds to days; roll calendars.
\sfield{Backtest honestly} The implied prices the engine would have generated; roll-period volumes.
\sfield{Sources} Book 1, chapter 19; this chapter.
\end{strategyfile}

\begin{strategyfile}{Implied-price arbitrage between spread and outright books}\label{strat:hf:futures-market-making:implied}
\sfield{Who pays you, and why} Spread-book quoters whose prices lag the outrights.
\sfield{Instruments and venues} A spread book and its two outright books, where the exchange's implied matching does not close every link.
\sfield{Signal} The direct spread price outside the implied range from the outrights.
\sfield{Sizing and execution} Take the spread and the two outright legs at once; size to the smallest of the three top-of-book quantities.
\sfield{Costs} Three fees; \omterm{def:hf:index-and-basket-arbitrage:legging}{legging risk} if one leg misses (chapter 13).
\sfield{How it dies} The exchange's implied engine, and faster traders: in the simulation the gaps appear only when the spread books lag by several
events.
\sfield{Horizon, capacity, infrastructure} Microseconds; the three books' data on one clock.
\sfield{Backtest honestly} Sequenced book events across the three books; whether the implied engine would have matched first.
\sfield{Sources} This chapter.
\end{strategyfile}

\begin{strategyfile}{Short-rate futures strip making}\label{strat:hf:futures-market-making:strip}
\sfield{Who pays you, and why} Rates hedgers and speculators trading any point of the strip, and packs and bundles.
\sfield{Instruments and venues} A strip of quarterly short-rate futures, its spreads, butterflies, packs and bundles (One Quant Book 2, chapter 8).
\sfield{Signal} A curve model: level, slope, curvature, central-bank meeting dates, fitted to all the strip's books at once.
\sfield{Sizing and execution} Quote every book from the model; hedge any fill in the cheapest instrument that removes its curve exposure.
\sfield{Costs} Fees, \omterm{def:hf:futures-market-making:over}{over-quoting}'s inventory risk in pro-rata books, model error around central-bank meetings.
\sfield{How it dies} Central-bank surprises move the curve in ways the model's factors do not span.
\sfield{Horizon, capacity, infrastructure} Seconds to days; hundreds of books from one model.
\sfield{Backtest honestly} The meeting calendar as known at the time; all books' allocation rules.
\sfield{Sources} One Quant Book 2, chapter 8.
\end{strategyfile}

\begin{strategyfile}{Inter-commodity spread trading}\label{strat:hf:futures-market-making:intercommodity}
\sfield{Who pays you, and why} Hedgers of processing margins and curve positions who trade the spread as one instrument.
\sfield{Instruments and venues} Two related futures and, where listed, their inter-commodity spread book.
\sfield{Signal} The spread's deviation from a model of the two prices' relation.
\sfield{Sizing and execution} Quote the spread; hedge fills in the leg with the better liquidity.
\sfield{Costs} Two legs' fees, margin (reduced by the exchange's spread credit), the relation's drift.
\sfield{How it dies} Structural change in the relation (a refinery outage, a new delivery rule).
\sfield{Horizon, capacity, infrastructure} Minutes to weeks.
\sfield{Backtest honestly} Margin credits as they were; delivery and contract changes by date.
\sfield{Sources} Book 1, chapter 20.
\end{strategyfile}

\section{Tutorial: showing ten times what you want}\label{tut:hf:futures-market-making:tut}

\begin{tutorial}
\textbf{Goal.} Find the equilibrium size makers show in a pro-rata book, its cost to them, and the implied arbitrages a lagging spread book leaves.
\textbf{End state:} \cref{fig:hf:futures-market-making:over,fig:hf:futures-market-making:implied}.
\begin{tutsteps}
\item \textbf{Allocation} of one aggressor under three rules with \texttt{firm.match} (\texttt{hf\_futures.allocation\_example}).
\item \textbf{The game}: best responses to the others' total size, damped, until they settle.
\omcode{code/firm/futmm/firm_futmm.py}{45}{62}{A maker's best size against the others' total; the symmetric equilibrium by damped best responses.}\label{lst:hf:futures-market-making:game}
\item \textbf{Stress}: the fills a 99.9th-percentile aggressor gives each maker.
\item \textbf{The strip}: four quarterly futures, three spread books quoting a lagged difference; scan for implied arbitrages.
\omcode{code/firm/futmm/firm_futmm.py}{99}{114}{Direct spread quotes against the spread implied by the outrights, with \texttt{firm.match.implied\_in}.}\label{lst:hf:futures-market-making:scan}
\end{tutsteps}
\textbf{What to change next.} Add a top-order allocation of 40\% to the game and see who gets it; let one maker predict large aggressors and cancel
before them; run the pro-rata book on \texttt{firm.exchsim}'s \texttt{pro\_rata\_futures} preset.
\end{tutorial}

\begin{center}
\small
\begin{tabular}{@{}lrrrrrr@{}}
\toprule
makers & 2 & 5 & 10 & 12 & 15 & 20\\
\midrule
size shown (lots) & 110 & 143 & 230 & 298 & 619 & 5,000\\
\omterm{def:hf:futures-market-making:over}{over-quoting} factor & 1.1 & 1.4 & 2.3 & 3.0 & 6.2 & 50\\
utility, equilibrium & 41.2 & 29.6 & 19.5 & 16.9 & 13.9 & 12.3\\
utility, all showing 100 & 41.4 & 31.2 & 22.3 & 20.0 & 17.4 & 14.3\\
fill from a 7,907-lot order & 110 & 143 & 230 & 298 & 527 & 395\\
depth / mean order & 0.5 & 1.8 & 5.6 & 8.8 & 23 & 245\\
\bottomrule
\end{tabular}
\end{center}

\section{Build: the futures market-making module}\label{bld:hf:futures-market-making:futmm}

\begin{build}
\bfield{Purpose} Model pro-rata fills and the \omterm{def:hf:futures-market-making:over}{over-quoting} equilibrium, and scan a strip for \omterm{def:hf:futures-market-making:implied}{implied-price arbitrages}.
\bfield{Interface} \texttt{fills(s, X)}, \texttt{utility}, \texttt{best\_response}, \texttt{equilibrium(n, X, h, gamma, grid)},
\texttt{fifo\_size}, \texttt{aggressors}, \texttt{Strip(n\_contracts, steps, seed, lag)} with \texttt{outright(t, i)} and \texttt{spread(t, i)},
\texttt{scan(strip)}. Built on \texttt{firm.match} (Book 1): \texttt{pro\_rata}, \texttt{configurable}, \texttt{implied\_in}.
\bfield{Rules} Fills are continuous shares of $\min(X,S)$; the grid's top is the exchange's maximum order size.
\bfield{Acceptance tests} \texttt{code/firm/futmm/tests/}: continuous fills within rounding of \texttt{firm.match.pro\_rata}; alone, the best response
is the time-priority size; the best response grows with the crowd; the equilibrium grows with the number of makers and leaves them worse off than
showing what they want; no arbitrage without lag, more with more lag.
\bfield{Stretch} Top-order and lead-market-maker slices in the game; the game on \texttt{firm.exchsim}'s pro-rata preset; packs and butterflies in
the scanner.
\end{build}

\begin{omsources}
\item J.~Field, J.~Large, Pro-rata matching and one-tick futures markets, CFS Working Paper 2008/40, 2008.
\item Databento, CME matching algorithms explained, 1 August 2025.
\end{omsources}

\section{Exercises}

\begin{exercise}[$\star$]\label{exo:hf:futures-market-making:1}
Four orders of 100 (the top order), 400, 300 and 200 lots rest at a price. Allocate an aggressor of 250 under time priority, pure pro rata, and top
order then pro rata.
\end{exercise}

\begin{exercise}[$\star$]\label{exo:hf:futures-market-making:2}
Ten makers each show 230 lots. What does an aggressor of 1,000 lots give each, and one of 7,907?
\end{exercise}

\begin{exercise}[$\star$]\label{exo:hf:futures-market-making:3}
The front contract is bid 9,612 and offered 9,613, the next bid 9,598 and offered 9,599. What is the implied spread market? A spread book is offered at
12: what do you do?
\end{exercise}

\begin{exercise}[$\star\star$]\label{exo:hf:futures-market-making:4}
Why is the equilibrium a prisoner's dilemma?
\end{exercise}

\begin{exercise}[$\star\star$]\label{exo:hf:futures-market-making:5}
Why do pro-rata books have cancellation rates above 96\%?
\end{exercise}

\begin{exercise}[$\star\star$]\label{exo:hf:futures-market-making:6}
How does a top-order allocation change a maker's incentives?
\end{exercise}

\begin{exercise}[$\star\star\star$]\label{exo:hf:futures-market-making:7}
\emph{Coding.} Double the \omterm{def:hf:inventory-models:penalty}{inventory penalty} (\texttt{GAMMA = 0.02}). What size would a maker show under time priority, and what do ten makers show?
\end{exercise}

\begin{exercise}[$\star\star\star$]\label{exo:hf:futures-market-making:8}
\emph{Find the flaw.} ``Depth at the best is 20,000 lots, so our 500-lot order will not move the price and we will get 500 lots.''
\end{exercise}

\section{Problem: Showing Ten Times What You Want}

\begin{problem}[{Weekend problem --- showing ten times what you want}]\label{pb:hf:futures-market-making:1}
Market makers in a pro-rata short-rate futures book choose how much to show, and a spread book beside it lags.

\medskip\noindent\textbf{Part I --- Allocation.}
\begin{enumerate}
\item Describe time priority, pro rata and top-order allocation, and the dated box's description of one exchange.
\item Allocate the example aggressor under each rule.
\item Define \omterm{def:hf:futures-market-making:over}{over-quoting}.
\item What did Field and Large observe in pro-rata markets?
\end{enumerate}

\medskip\noindent\textbf{Part II --- The game.}
\begin{enumerate}[resume]
\item Write a maker's fill and utility.
\item What size does a maker show under time priority, and why?
\item How is the symmetric equilibrium found?
\item Why does the spiral stop only at the cap with twenty makers?
\end{enumerate}

\medskip\noindent\textbf{Part III --- Measurements.}
\begin{enumerate}[resume]
\item Give the \omterm{def:hf:futures-market-making:over}{over-quoting} factor for 2, 10, 15 and 20 makers.
\item Compare equilibrium utility with everyone showing 100.
\item What does the 99.9th-percentile aggressor do to each maker?
\item Define \omterm{def:hf:futures-market-making:implied}{implied-price arbitrage} and give the counts by lag.
\end{enumerate}

\medskip\noindent\textbf{Part IV --- The verdict.}
\begin{enumerate}[resume]
\item State the \emph{named result}: the equilibrium \omterm{def:hf:futures-market-making:over}{over-quoting} factor in a pro-rata book, and the inventory risk it creates when an unexpectedly
large order fills everyone at once.
\item How should a maker size its display, knowing the equilibrium?
\item What could an exchange change to stop the spiral?
\item Why do strip makers quote hundreds of books from one model?
\item When does an inter-commodity spread stop being a hedge?
\item Which of the five strategy files depends most on speed?
\item How would you backtest a pro-rata strategy when others' sizes are unobservable?
\item In one sentence: what does size buy in a pro-rata book?
\end{enumerate}
\end{problem}

\section{Interview questions}

\begin{interviewq}[$\star$ \iqroles{trader}]\label{iq:hf:futures-market-making:1}
Why is the bid--ask spread of a short-rate future almost always one tick?
\end{interviewq}

\begin{interviewq}[$\star\star$ \iqroles{researcher}]\label{iq:hf:futures-market-making:2}
Derive the best response in the pro-rata game when aggressors never exhaust the book.
\end{interviewq}

\begin{interviewq}[$\star\star$ \iqroles{trader}]\label{iq:hf:futures-market-making:3}
You show 2,000 lots and a 10,000-lot order hits the bid. What happens next, and what should your system do?
\end{interviewq}

\begin{interviewq}[$\star\star$ \iqroles{developer}]\label{iq:hf:futures-market-making:4}
Implement pro-rata allocation with rounding down, a minimum allocation and FIFO leftovers. What are the edge cases?
\end{interviewq}

\begin{interviewq}[$\star\star$ \iqroles{risk}]\label{iq:hf:futures-market-making:5}
Your firm's displayed size in a pro-rata book is fifty times its position limit. Is that a breach?
\end{interviewq}

\begin{interviewq}[$\star\star\star$ \iqroles{researcher}]\label{iq:hf:futures-market-making:6}
How would you estimate, from public market data, how much of the displayed depth in a pro-rata book is \omterm{def:hf:futures-market-making:over}{over-quoting}?
\end{interviewq}
